% ================================================================================
% PECA 131 / PAPEL DE TRABALHO PT - ANALISE DE CAPACIDADE CIVIL NO CAF
% ================================================================================
% Data: 03/dezembro/2025
% Objetivo: Documentar inconsistencias de capacidade civil identificadas no CAF
%           atraves de cruzamento com bases SISOBI (obitos) e RFB (CPF) via LABCONTAS.
% Referencia: peca131_pt_capacidade_civil.pdf
% ================================================================================

\documentclass[11pt,a4paper]{article}

% ========== PACOTES PADRAO ==========
\usepackage[utf8]{inputenc}
\usepackage[brazil]{babel}
\usepackage[T1]{fontenc}
\usepackage{lmodern}
\usepackage[margin=2.5cm]{geometry}
\usepackage{graphicx}
\usepackage{float}
\usepackage{longtable}
\usepackage{booktabs}
\usepackage{xcolor}
\usepackage{colortbl}
\usepackage{fancyhdr}
\usepackage{amsmath, amssymb}
\usepackage{listings}
\usepackage{enumitem}
\usepackage[hidelinks]{hyperref}
\usepackage[breakable]{tcolorbox}

% ========== CORES TCU OFICIAIS ==========
\definecolor{tcured}{RGB}{204,0,0}      % Vermelho TCU #CC0000
\definecolor{tcublue}{RGB}{0,51,102}    % Azul TCU #003366
\definecolor{tcugreen}{RGB}{0,153,76}   % Verde TCU #006633

% ========== CONFIGURACAO DE CABECALHO ==========
\pagestyle{fancy}
\fancyhead[L]{Peca 131 / PT - Capacidade Civil no CAF}
\fancyhead[R]{TCU}
\setlength{\headheight}{14pt}

% ========== CONFIGURACAO DE LISTINGS ==========
\lstset{
    basicstyle=\ttfamily\small,
    frame=single,
    breaklines=true,
    numbers=left,
    numberstyle=\tiny\color{gray},
    keywordstyle=\color{tcublue}\bfseries,
    commentstyle=\color{tcugreen}\itshape,
    stringstyle=\color{tcured},
    columns=fullflexible
}

% ========== INFORMACOES DO DOCUMENTO ==========
\title{Papel de Trabalho de Auditoria Tecnica\\
\textbf{Peca 131: Analise de Capacidade Civil dos Responsaveis no CAF}}
\author{Equipe de Auditoria --- Fiscalizacao CAF-MAPA}
\date{03 de dezembro de 2025}

\begin{document}

\maketitle

% ================================================================================
\section{Objetivo}
% ================================================================================

Este papel de trabalho documenta a \textbf{analise de capacidade civil dos responsaveis}
cadastrados no Sistema CAF, verificando:

\begin{itemize}
    \item Responsaveis falecidos ainda ativos no cadastro (cruzamento com SISOBI);
    \item Menores de 16 anos cadastrados como responsaveis (absolutamente incapazes);
    \item Adolescentes de 16-17 anos nao emancipados (relativamente incapazes);
    \item Divergencias de data de nascimento entre CAF e Receita Federal;
    \item CPFs nao localizados na base da Receita Federal.
\end{itemize}

\textbf{Fundamentacao Legal}: Artigos 3\textordmasculine{} e 4\textordmasculine{} do Codigo Civil
(Lei 10.406/2002) definem a incapacidade civil, tornando nulos ou anulaveis os atos
praticados por incapazes.

% ================================================================================
\section{Metodologia}
% ================================================================================

\subsection{Base de Dados}

A analise foi realizada no ambiente \textbf{LABCONTAS} do TCU, utilizando:

\begin{itemize}
    \item \textbf{Fonte CAF}: \texttt{[BD\_AGRO\_CAF\_AUD\_2025].[dbo].[S\_*]}
    \item \textbf{Fonte Obitos}: \texttt{[BD\_SISOBI].[dbo].[SISOBI]}
    \item \textbf{Fonte CPF}: \texttt{[BD\_RECEITA].[dbo].[CPF]}
    \item \textbf{Populacao}: Responsaveis de Unidades Familiares ativas
\end{itemize}

\subsection{Criterio de Selecao}

Foi aplicado o \textbf{JOIN Triplo canonico} para identificar responsaveis de UFs ativas:

\begin{lstlisting}[language=SQL, caption={Selecao de responsaveis ativos}]
SELECT DISTINCT
    pf.nr_cpf,
    pf.dt_nascimento AS dt_nascimento_caf,
    uf.dt_ativacao,
    uf.dt_atualizacao
FROM [BD_AGRO_CAF_AUD_2025].[dbo].[S_RESPONSAVEL_UFPR] r
INNER JOIN [BD_AGRO_CAF_AUD_2025].[dbo].[S_UNIDADE_FAMILIAR] uf
    ON r.id_unidade_familiar = uf.id_unidade_familiar
INNER JOIN [BD_AGRO_CAF_AUD_2025].[dbo].[S_PESSOA_FISICA] pf
    ON r.id_pessoa_fisica = pf.id_pessoa_fisica
WHERE uf.id_tipo_situacao_unidade_familiar = 1  -- UF ativa
  AND r.st_excluido = 0;
\end{lstlisting}

\subsection{Cruzamentos Externos}

\begin{enumerate}
    \item \textbf{SISOBI}: Verificacao de obito pelo CPF do responsavel;
    \item \textbf{RFB}: Validacao de data de nascimento e existencia do CPF.
\end{enumerate}

% ================================================================================
\section{Resultados}
% ================================================================================

\subsection{Responsaveis Falecidos (SISOBI)}

O cruzamento com a base SISOBI identificou \textbf{12.820 responsaveis falecidos}
ainda ativos no CAF.

\begin{table}[H]
\centering
\caption{Obitos por Periodo}
\label{tab:obitos}
\begin{tabular}{lrrl}
\toprule
\textbf{Periodo do Obito} & \textbf{Quantidade} & \textbf{\%} & \textbf{Observacao} \\
\midrule
Antes de 2020 & 132 & 1,0\% & Obitos antigos \\
2020-2022 & 2.891 & 22,5\% & Periodo COVID \\
2023 & 3.647 & 28,4\% & --- \\
2024 & 4.812 & 37,5\% & --- \\
2025 & 1.338 & 10,4\% & Ano corrente \\
\midrule
\textbf{Total} & \textbf{12.820} & \textbf{100\%} & --- \\
\bottomrule
\end{tabular}
\end{table}

\subsubsection{Casos Extremos de Obitos Antigos}

\begin{table}[H]
\centering
\caption{Obitos Mais Antigos com Cadastros Ainda Ativos}
\label{tab:obitos_antigos}
\begin{tabular}{cccr}
\toprule
\textbf{Ano Obito} & \textbf{Data Ativacao} & \textbf{Data Atualizacao} & \textbf{Anos Pos-Obito} \\
\midrule
\textbf{1985} & 28/ago/2024 & 08/out/2024 & \textbf{39 anos} \\
1997 & --- & --- & 28 anos \\
1998 & --- & --- & 27 anos \\
2011 & 05/mar/2025 & 05/mar/2025 & 14 anos \\
2011 & 05/mar/2025 & 05/mar/2025 & 14 anos \\
\bottomrule
\end{tabular}
\end{table}

\textbf{Achado critico}: Responsavel falecido ha \textbf{39 anos} (1985) com cadastro
\textbf{ativado em agosto/2024} e \textbf{atualizado em outubro/2024}.

% ----------------

\subsection{Menores Absolutamente Incapazes}

Foram identificados \textbf{40 menores de 16 anos} cadastrados como responsaveis de
Unidades Familiares --- atos \textbf{nulos} conforme Art. 3\textordmasculine{} do Codigo Civil.

\begin{table}[H]
\centering
\caption{Distribuicao por Idade na Ativacao}
\label{tab:menores}
\begin{tabular}{crr}
\toprule
\textbf{Idade na Ativacao} & \textbf{Quantidade} & \textbf{Observacao} \\
\midrule
-1 ano (erro de data) & 5 & Divergencia com RFB \\
0 ano (recem-nascido) & 9 & Impossivel \\
1 ano & 10 & Impossivel \\
2-6 anos & 1 & Criancas \\
7-10 anos & 3 & Criancas \\
\textbf{11 anos} & \textbf{1} & \textbf{Caso extremo} \\
13-15 anos & 11 & Pre-adolescentes \\
\midrule
\textbf{Total} & \textbf{40} & --- \\
\bottomrule
\end{tabular}
\end{table}

\textbf{Achado critico}: \textbf{Crianca de 11 anos} cadastrada como responsavel de
Unidade Familiar --- caso de absoluta incapacidade civil.

% ----------------

\subsection{Adolescentes Nao Emancipados}

Foram identificados \textbf{210 adolescentes de 16-17 anos} nao emancipados cadastrados
como responsaveis --- atos \textbf{anulaveis} conforme Art. 4\textordmasculine{}, I do Codigo Civil.

\begin{table}[H]
\centering
\caption{Adolescentes 16-17 Anos}
\label{tab:adolescentes}
\begin{tabular}{lrl}
\toprule
\textbf{Categoria} & \textbf{Quantidade} & \textbf{Status} \\
\midrule
16 anos & 98 & Nao emancipados \\
17 anos & 112 & Nao emancipados \\
\midrule
\textbf{Total} & \textbf{210} & \texttt{st\_declarante\_emancipado = False} \\
\bottomrule
\end{tabular}
\end{table}

\textbf{Observacao}: O campo \texttt{st\_declarante\_emancipado = False} indica expressamente
que estes adolescentes \textbf{nao sao emancipados}, configurando incapacidade relativa.

% ----------------

\subsection{Divergencias de Data de Nascimento}

O cruzamento CAF vs RFB identificou \textbf{12.777 divergencias} de data de nascimento.

\begin{table}[H]
\centering
\caption{TOP 10 Maiores Divergencias}
\label{tab:divergencias}
\begin{tabular}{cccr}
\toprule
\textbf{Data CAF} & \textbf{Data RFB} & \textbf{Diferenca} & \textbf{Observacao} \\
\midrule
10/nov/1879 & 10/nov/1989 & \textbf{110 anos} & Seculo errado \\
10/dez/2028 & 20/nov/1921 & 107 anos & Data futura no CAF \\
12/out/2082 & 12/out/1982 & 100 anos & Data futura no CAF \\
15/mar/2045 & 15/mar/1945 & 100 anos & Data futura no CAF \\
04/out/2049 & 04/out/1949 & 100 anos & Data futura no CAF \\
21/jul/1885 & 21/jul/1985 & 100 anos & Seculo errado \\
20/jan/2039 & 20/jan/1939 & 100 anos & Data futura no CAF \\
01/jan/2019 & 01/jan/1919 & 100 anos & Data futura no CAF \\
18/set/2024 & 18/set/1924 & 100 anos & Data futura no CAF \\
10/abr/2045 & 10/abr/1945 & 100 anos & Data futura no CAF \\
\bottomrule
\end{tabular}
\end{table}

\textbf{Achado critico}: Divergencia maxima de \textbf{110 anos} entre data no CAF
(10/nov/1879) e data na RFB (10/nov/1989).

% ----------------

\subsection{Datas de Nascimento Futuras}

Foram identificados \textbf{48 registros} com data de nascimento no \textbf{futuro}
(apos 2025), resultando em ``idades negativas''.

\begin{table}[H]
\centering
\caption{Datas de Nascimento Futuras por Ano}
\label{tab:futuras}
\begin{tabular}{crl}
\toprule
\textbf{Ano no CAF} & \textbf{Quantidade} & \textbf{Observacao} \\
\midrule
2084 & 1 & 59 anos no futuro \\
2083 & 1 & 58 anos no futuro \\
2082 & 1 & 57 anos no futuro \\
2049 & 8 & --- \\
2048 & 5 & --- \\
2045-2047 & 9 & --- \\
2028-2044 & 23 & --- \\
\midrule
\textbf{Total} & \textbf{48} & Todas invalidas \\
\bottomrule
\end{tabular}
\end{table}

\textbf{Observacao}: Estes registros evidenciam \textbf{falha critica de validacao} ---
o sistema aceita datas de nascimento no futuro sem qualquer verificacao.

% ----------------

\subsection{CPFs Nao Localizados na RFB}

Foram identificados \textbf{520 CPFs} cadastrados no CAF que \textbf{nao existem}
na base da Receita Federal.

% ================================================================================
\section{Analise}
% ================================================================================

\subsection{Consolidacao dos Achados}

\begin{table}[H]
\centering
\caption{Resumo das Inconsistencias de Capacidade Civil}
\label{tab:resumo}
\begin{tabular}{lrrr}
\toprule
\textbf{Categoria} & \textbf{Quantidade} & \textbf{\%} & \textbf{Impacto Juridico} \\
\midrule
Responsaveis falecidos & 12.820 & 49,1\% & Cadastro invalido \\
Menores $<$16 anos & 40 & 0,2\% & Atos nulos \\
Adolescentes 16-17 & 210 & 0,8\% & Atos anulaveis \\
Divergencias nascimento & 12.777 & 48,9\% & Dado incorreto \\
Datas futuras & 48 & 0,2\% & Validacao ausente \\
CPFs inexistentes & 520 & 2,0\% & Cadastro irregular \\
\midrule
\textbf{Total de registros} & \textbf{26.415} & --- & --- \\
\bottomrule
\end{tabular}
\end{table}

\textbf{Nota}: Ha sobreposicao entre categorias (ex: divergencias podem incluir datas futuras).
O total consolidado de registros unicos com problemas e de aproximadamente \textbf{26.000 casos}.

\subsection{Falhas de Validacao Identificadas}

\begin{enumerate}
    \item \textbf{Ausencia de cruzamento com SISOBI}: O sistema nao verifica obito do
    responsavel antes de permitir ativacao ou atualizacao de cadastro;

    \item \textbf{Nao verificacao de maioridade}: O sistema permite cadastrar menores de
    16 anos como responsaveis, sem alertas ou bloqueios;

    \item \textbf{Campo de emancipacao nao validado}: Adolescentes 16-17 anos nao emancipados
    (\texttt{st\_declarante\_emancipado = False}) sao aceitos como responsaveis;

    \item \textbf{Validacao de data insuficiente}: O sistema aceita datas de nascimento no
    futuro e datas incompativeis com a idade do titular do CPF.
\end{enumerate}

% ================================================================================
\section{Conclusao}
% ================================================================================

A analise de capacidade civil dos responsaveis no CAF revela \textbf{falhas sistemicas
de validacao}:

\begin{enumerate}
    \item \textbf{12.820 responsaveis falecidos} ainda ativos no cadastro, incluindo
    caso extremo de obito em \textbf{1985} (39 anos atras) com atualizacao em \textbf{2024};

    \item \textbf{40 menores absolutamente incapazes} ($<$16 anos), incluindo caso de
    \textbf{crianca de 11 anos} cadastrada como responsavel;

    \item \textbf{210 adolescentes relativamente incapazes} (16-17 anos nao emancipados);

    \item \textbf{12.777 divergencias} de data de nascimento entre CAF e RFB, incluindo
    divergencia maxima de \textbf{110 anos};

    \item \textbf{48 datas de nascimento no futuro}, evidenciando ausencia total de
    validacao temporal.
\end{enumerate}

\textbf{Impacto juridico}: Os atos praticados por absolutamente incapazes sao \textbf{nulos}
(Art. 166, I do CC), enquanto os praticados por relativamente incapazes sao \textbf{anulaveis}
(Art. 171, I do CC). A manutencao de cadastros de pessoas falecidas pode caracterizar
irregularidade administrativa e, em casos extremos, fraude.

% ================================================================================
\section{Anexos}
% ================================================================================

\subsection{Parametros de Conexao LABCONTAS}

\begin{verbatim}
Ambiente: LABCONTAS (SQL Server)
Base CAF: [BD_AGRO_CAF_AUD_2025]
Base SISOBI: [BD_SISOBI]
Base RFB: [BD_RECEITA]
Data de execucao: Novembro/2025
\end{verbatim}

\subsection{Artefatos Gerados (CSVs)}

Os seguintes arquivos CSV foram exportados das queries executadas no LABCONTAS e
servem como evidencia primaria dos achados:

\begin{itemize}
    \item \texttt{ACHADO\_CRITICO\_responsaveis\_com\_obito.csv} (12.820 registros)
    \item \texttt{ACHADO\_CRITICO\_menores\_16\_anos.csv} (40 registros)
    \item \texttt{ACHADO\_16\_17\_anos\_nao\_emancipados.csv} (210 registros)
    \item \texttt{ACHADO\_datas\_nascimento\_divergentes.csv} (12.777 registros)
    \item \texttt{ACHADO\_cpf\_nao\_encontrado\_receita.csv} (520 registros)
\end{itemize}

\subsection{Artefatos deste Documento}

\begin{itemize}
    \item \texttt{peca131\_pt\_capacidade\_civil.tex} --- Este documento (codigo fonte)
    \item \texttt{peca131\_pt\_capacidade\_civil.pdf} --- Versao compilada (Peca 131)
\end{itemize}

\end{document}
